% !TEX root = ../main.tex

\chapter{全文总结与展望}

\section{全文总结}

本文围绕金融大规模数据异常检测系统的实现与优化问题展开研究，针对传统异常检测系统在大规模场景下面临的异常评分复杂度高、候选召回吞吐不足、CPU 与 GPU 之间数据搬运频繁以及多 GPU 扩展效率受限等问题，提出并实现了一套基于并行计算的 PyTOD--FlashANNS 融合异常检测系统。全文研究重点不在于提出新的异常检测理论模型，而在于将候选召回、异常评分、数据路径优化与多 GPU 扩展统一纳入同一系统框架，并验证其在金融近实时异常检测场景中的工程可行性\citep{zhao2022tod,xiao2025flashanns,cuvs_multi_gpu_nn}。

首先，在系统总体架构方面，本文提出了以 PyTOD 为异常评分主干、以 FlashANNS 为候选召回模块的两阶段融合框架。本文将二者统一组织到一条候选召回---异常评分执行链路中，使系统能够先在大规模样本中快速召回候选集合，再在候选集合上执行张量化异常评分，从而在保证检测效果可接受的前提下，显著降低全量异常评分所带来的时间与空间开销\citep{zhao2022tod,xiao2025flashanns}。

其次，在数据准备与任务建模方面，本文从金融交易数据、账户行为数据和图金融数据出发，对输入形式、向量化表示、候选召回目标和异常评分目标进行了统一建模，并形成了``真实/半真实数据用于有效性验证，高保真合成数据用于大规模系统实验''的数据组织方案。通过引入 IEEE-CIS Fraud Detection、Credit Card Fraud Detection、DGraphFin 以及 IBM AML/AMLSim 等数据来源，本文构建了 1M、10M 和 100M 三个规模层级的实验体系，使论文能够同时覆盖方法可用性验证和系统吞吐扩展验证两类目标\citep{ibm_aml_data,ibm_amlsim,ieee_cis_fraud,dalpozzolo2015creditcard,dgraphfin_baseline}。

再次，在单 GPU 场景下，本文提出了 GPU 主导的端到端执行链路，将查询特征、候选结果与异常评分中间张量尽量保持在 GPU 侧，仅在必要边界通过 pinned memory 与主机交互，从而压缩系统中的非必要数据传输和阶段同步开销\citep{xiao2025flashanns,faiss_gpu_wiki,rapids_rmm_pinned}。与此同时，本文进一步利用 RAPIDS 对表格型数据处理和中间结果组织进行后端重构，并借助 FAISS-GPU 对候选检索后端进行工程增强，使系统在单卡场景下从``模块可运行''演进为``链路低开销、执行更连续''的高吞吐实现。

进一步地，在多 GPU 场景下，本文以吞吐优先为目标，提出了数据并行与索引复制的主扩展方案，使每张 GPU 均持有完整索引并独立完成本地候选召回与异常评分，CPU 主要承担控制面调度与最终小规模结果汇总职责\citep{cuvs_multi_gpu_nn,cuvs_dist_mode}。在此基础上，本文讨论了索引复制不可行时的兜底路径：通过 GPU 侧局部结果组织、NCCL 设备间通信和 GPU 侧最终归并，使 CPU 尽量继续退出数据平面\citep{nccl_stream_semantics}。同时，本文引入拓扑感知的 I/O lane 绑定与准入控制机制，以缓解多 GPU 条件下的路径冲突和尾延迟放大问题\citep{gds_overview,gds_design_guide}。

最后，在综合实验与结果分析方面，本文从异常检测效果、候选召回质量、单卡端到端吞吐、多 GPU 扩展能力、CPU 占用、GPU 利用率、I/O 路径占用和延迟分位数等多个维度对系统进行了统一评估。实验结果表明，融合架构在 1M 规模真实/半真实数据上能够保持合理的检测效果；在 10M 和 100M 规模高保真合成数据上，单卡 GPU 主导执行链路能够显著减少 CPU$\leftrightarrow$GPU 往返和 I/O 等待，而多 GPU 数据并行与索引复制方案能够较稳定地支撑系统吞吐扩展\citep{xiao2025flashanns,cuvs_multi_gpu_nn,ibm_aml_data}。

综合来看，本文的主要研究结论可以归纳为三点：
\begin{enumerate}
  \item PyTOD 与 FlashANNS 的融合能够有效打通高吞吐候选召回与高效异常评分之间的链路，使系统从全量高复杂度执行转向候选缩减后的精排评分路径。
  \item 单 GPU 场景下的系统瓶颈主要来自数据路径与阶段同步，而不是单个算子本身；通过 GPU 主导端到端执行链路、RAPIDS 后端重构和 FAISS-GPU 工程增强，可以显著提升端到端性能。
  \item 多 GPU 场景下更适合采用数据并行与索引复制作为吞吐优先主方案，只有在索引无法复制时才进入 GPU 侧跨卡归并与设备间通信路径；同时，多 GPU 上限与存储拓扑和 I/O 路径组织密切相关，因此拓扑感知 I/O 控制应被视为系统设计的重要组成部分。
\end{enumerate}

表~\ref{tab:overall-summary} 对全文主要工作与章节贡献进行了对应汇总，图~\ref{fig:overall-roadmap} 则从技术路线和结论闭环两个角度概括了本文的整体逻辑。

\paragraph{图表预留说明：}
\begin{itemize}
  \item \textbf{表 7-1 全文主要工作与对应章节总结表。} 将全文的主要研究工作、核心章节和实验验证结果进行一一对应。
  \item \textbf{图 7-1 全文技术路线与结论闭环图。} 展示从问题提出、融合系统构建、单卡优化、多 GPU 扩展到综合实验验证的完整闭环。
\end{itemize}

\begin{table}[htbp]
  \caption{全文主要工作与对应章节总结}
  \label{tab:overall-summary}
  \centering
  \thesisplaceholderbox[0.95\linewidth]{正式版补齐表~7-1}
\end{table}

\begin{figure}[htbp]
  \centering
  \thesisplaceholderbox{正式版补齐图~7-1}
  \caption{全文技术路线与结论闭环图}
  \label{fig:overall-roadmap}
\end{figure}

\section{未来工作展望}

尽管本文围绕金融大规模异常检测系统的实现与优化进行了较为系统的研究，并在 PyTOD--FlashANNS 融合架构、单卡数据路径压缩和多 GPU 扩展方面取得了一定成果，但从更长期的工程部署与学术研究角度看，仍有若干值得进一步深入的方向。

\subsection{面向更真实金融场景的数据验证}

未来工作需要进一步争取接入更接近真实生产环境的金融业务数据，在严格遵守隐私保护与合规要求的前提下，验证系统在更复杂分布、更强业务噪声和更显著概念漂移条件下的表现。

\subsection{面向多机场景的分布式扩展}

未来可考虑将本文提出的``控制面与数据平面分离''思想扩展到多节点环境，在更大规模集群中研究查询分发、跨节点索引一致性、远程结果归并和跨节点通信调度问题。

\subsection{面向更深层存储路径的优化}

未来可以进一步引入 GPUDirect Storage、RDMA 等技术，探索更少主机参与、更低拷贝代价的存储访问路径，并在查询调度与存储访问之间建立更细粒度的协同策略。

\subsection{面向图异常检测与时序异常检测的统一扩展}

未来可在保持系统总体架构不变的前提下，引入图嵌入、图神经网络、时序编码或基于事件流的特征聚合机制，使图异常检测与时序异常检测在不牺牲系统吞吐的前提下融入现有候选召回---异常评分框架。

\subsection{面向国产 GPU 平台的验证与部署}

未来还应关注本文方案在国产 GPU 平台上的验证与部署问题，围绕算子支持、向量检索实现、设备间通信机制与存储接口适配，研究如何构建等价或近似等价的工程实现路径。
